%%%%%%%%%%%%%%%%%%%%%%%%%%%%%%%%%
%
%       EXERCÍCIO
%
%%%%%%%%%%%%%%%%%%%%%%%%%%%%%%%%%

\definevalorquestao

\begin{exercicioBanco}[\valorquestao]
Os valores de \(x \in \mathbb{R}\) que satisfazem simultaneamente as condições \(\left(\dfrac{1}{5}\right)^{x^2} \le 5^{-4x}\) e \(x^2 \le 5\) são:

% Define as alternativas
\newcommand{\alternativas}{%
    \begin{center}
        \begin{tabularx}{\textwidth}{XXX}
            (a) \(x \le -\sqrt{5}\) ou \(x \ge \sqrt{5}\). &
            (b) \(-\sqrt{5} \le x \le \sqrt{5}\). &
            (c) \(0 \le x \le 4\). \\[5pt]
            (d) \(x \le 0\) ou \(x \ge 4\). &
            (e) \(-\sqrt{5} \le x \le 0\).
        \end{tabularx}
    \end{center}
}

% Define a resposta correta
\newcommand{\resposta}{E}

% Lógica condicional para exibição
\ifdefstring{\atividade}{lista}{%
    \alternativas
    \vspace{0.5em}
    
    \noindent\textbf{Resposta:} letra \textbf{\resposta}.
}{%
    \ifdefstring{\modo}{objetivo}{%
        \alternativas
    }{%
        % modo = discursiva → não mostra alternativas
    }
}
\end{exercicioBanco}

